\section{Relation to channel metrology and measurement information}
\label{sec:comparison84}

\subsection{Support spans, channel estimation and finite pairs}
For the measurement Kraus operators
$L_{j,a}=|j\rangle\langle a|\sqrt{E_j}$, the real Hermitian
Kraus-product span is exactly $\mathcal W_E$.  For the orbit
$L_{j,a}(t)=L_{j,a}e^{-itG}$ the effective generator is $G$.
Zhou--Jiang \cite[Theorems~1--3 and Section~V]{ZJ84} establish the
Kraus-span criterion for the linear or quadratic asymptotic channel
Fisher-information rate and give error-correction attainability.
We do not claim a new metrological dichotomy or a new correction
principle.  Theorems~\ref{thm:supportkernel84} and
\ref{thm:orbittrichotomy84} identify its complete support-kernel
form and give finite-angle trace-distance inequalities with an
explicit $4mt^2\|G\|_{\mathrm{op}}^2$ remainder.
Theorem~\ref{thm:curveclassification85} adds a normalized first-order
realization for general two-sided $C^2$ effect curves; its
$\sqrt N|t|+O(Nt^2)$ bound handles rank changes without a smooth
exact Kraus hypothesis.  Theorem~\ref{thm:controlstability85} separately
budgets certified control errors.  Neither constitutes a new claim
to the underlying metrological exponents.

Sieniawski--Demkowicz-Dobrza\'nski
\cite[Section~IV, equations~(16)--(23)]{SD76} integrate channel
Fisher-information bounds to bound adaptive discrimination and
compare metrological scaling with perfect finite discrimination.
The same horizontal method is a principal antecedent of
\eqref{eq:frontupper84}.  Our support-kernel identity is stated in
measurement coordinates and our lower \eqref{eq:finiteanglelower84}
comes from an actual corrected finite-angle channel.  It does not
claim the exact optimal tester or perfect discrimination.

\subsection{Fisher frames and covariance embeddings}
Saini--Kiukas--Burgarth--Gilchrist \cite[Theorem~1]{SKBG84} compare
classical and quantum Fisher information for a varying full-rank
\emph{state} measured by one fixed informationally complete POVM,
using a state-dependent frame operator.  Here the varying object
is the measurement, the tangent has coupled normalization, and the
loss permits $N$ adaptive uses with a reference.  In the rank-one
case our support span equals the effect span, so informational
completeness excludes the linear orbit regime by
Corollary~\ref{cor:orbitexamples84}.  That intersection does not
identify the two quadratic forms or their resource models.

Mayer--Yun \cite[Theorem~3.6, Definition~3.7 and Remark~4.4]{MY84}
use a kernel embedding of operator-valued measures, a kernel-dependent
discrepancy, and a tomographic Gram operator whose nullspace records
unobserved state directions.  Our support span contains entire
matrix blocks $P_j\mathcal H_dP_j$, rather than one observable
per effect; this distinction disappears for rank-one effects but
not in general.  Theorem~\ref{thm:supportkernel84} concerns a
coupled measurement tangent, and \eqref{eq:frontupper84} concerns
adaptive finite-use loss.  Neither kernel embedding nor injectivity
alone supplies those statements.

\subsection{Known symmetry and tomography primitives}
The full primary text of Yoshida--Okigami--Posta--Grinko
\cite{YOPG84} distinguishes known-symmetry estimation from known-pair
discrimination.  Theorem~7 gives invariant-state copy complexity
$\Theta((D_{\rm st}-1+\log(1/\eta))/\epsilon^2)$, where
$D_{\rm st}=\sum_\lambda m_\lambda^2$; $D_{\rm st}=1$ is a
known singleton.  Corollary~9 gives parallel covariant-channel bounds
\[
 O((D'_G+K_G\log(1/\eta))/\epsilon_\diamond^2),\qquad
 O((C'_G+\log(1/\eta))/\epsilon_{\rm tr}^2),
\]
in one-use halved diamond distance and normalized-Choi trace distance,
respectively, for $0<\epsilon\le1$, $0<\eta\le1/2$.
Here $D_G=\sum_\lambda m_\lambda M_\lambda$ counts Choi commutant
parameters before trace preservation,
$K_G=\sum_\lambda m_\lambda$,
$D'_G=D_G+K_G\log(4L)$, and
$C'_G=d_{\mathcal I}^{-1}(\sum_\lambda m_\lambda
\sqrt{d_\lambda M_\lambda})^2$; $L$ counts occupied blocks.
Theorem~13 does not give a general matching inverse-square diamond
lower: its bound is $\Omega_\xi(D_G/\epsilon_\diamond^{\xi/2})$
under its affine-dimension hypothesis.

For permutation covariance on $\ell$ systems of fixed local dimension
$q$, Theorems~12 and~15 match the Choi order
$\ell^{q^4-q^2}\epsilon_{\rm tr}^{-2}$ at constant confidence.
Theorem~18 approximates the Hayashi measurement in
$O((s+q)\operatorname{polylog}(s,q,\tau^{-1},\zeta^{-1}))$ gates,
with coupled output error at most $\tau$ except with probability
$\zeta$.  This is not a diamond-error certificate for our recovery.
Known-symmetry estimation, its parameter dimensions, and that
specialized gate construction do not themselves give our support
kernel or finite-pair curve theorem.  Nor do our theorems claim
faster learning on those symmetry classes.

For ordinary tomography the main-text substitutions remain explicit.
Mele--Bittel \cite[Theorem~III.3]{MB67} give the sufficient
constant-confidence order $O(dkr\epsilon^{-2})$ in halved diamond
loss, with $r\le dk$ for measurement channels.
Zambrano--Ramos-Calderer--Kueng \cite[Theorem~2]{ZRK76} give the
sufficient single-copy worst-input outcome-TV bound
\[
 m\ge\frac{8[d^2+\epsilon(d^2+1)/6]}{\epsilon^2}
          \log\frac{2^{k+1}9^{2d}}{\eta}.
\]
Either loss bounds the component norm error by $\epsilon$ after
output dephasing where needed.  A legal estimator need not be
balanced: apply Theorem~\ref{thm:affinerepair83} before using the
interior metric.  The resulting future loss is at most
$8k\sqrt N\epsilon$.  With $\epsilon=\delta/(8k\sqrt N)$ the
respective sufficient orders are
\[
 O(d^2k^4N\delta^{-2}),\qquad
 O\!\left(k^2N\delta^{-2}d^2[d+k+\log(1/\eta)]\right).
\]
These are upper-bound substitutions, not optimality claims about
the cited estimators.  Our component procedure credits the existing
binary operator-norm primitive and uses a different confidence
allocation.  The full calculations and confidence restrictions
are retained in the supplement, Section~\ref{sec:priority83}.

\subsection{What remains beyond the theorems}
The orbit theorem and Theorem~\ref{thm:curveclassification85} provide
matching scales on each fixed two-sided smooth curve with nonzero
first derivative, including second-order rank openings.  Their
constants need not be uniform across curves.  Arbitrary-pair
midpoint equivalence and full-body multi-outcome entropy are not
consequences of these local statements.  Growing-$k$ minimax sharpness,
polynomial entropy-optimal dictionaries, and general readout synthesis
remain separate mathematical questions.  Independent expert priority
assessment is not inferred from the absence of an identical statement
in this targeted comparison.

\subsection{Tangent neighborhoods and independent acquisition}
Theorem~\ref{thm:tubedichotomy86} fixes $E,H$ and an explicit
second-order remainder allowance, rather than assuming a particular
smooth family of Kraus operators.  It treats the full one-sided
positive tangent cone.  A nonzero missing-support compression has a
classical impossible-event witness; only within the cone's lineality
space does the support-generator test select the coherent branch.
The latter still uses the credited correction mechanism.
Theorem~\ref{thm:producttrichotomy86} separates these mechanisms by
an explicit product probe--reference restriction.  Its upper uses
standard root-fidelity inequalities and tensor multiplicativity
\cite[Chapter~3]{Watrous65}, not a new fidelity principle or a bound
on all nonadaptive entangled strategies.

The higher-order statement requires the finite remainder
\eqref{eq:powerjet86}.  Proposition~\ref{prop:mixedrates86} instead
proves a two-scale law when an intermediate support-opening term is
present.  These results address one-sided openings and selected
higher-order approaches without identifying a first jet with an
arbitrary-pair global geometry.  No general boundary entropy or
unrestricted growing-outcome learning law is inferred.
